\chapter{Implementation}
\label{ch:implementation}

\section{Introduction}

This chapter explains how the design of Chapter~4 was built. It describes the
technology stack, the ingestion pipeline that processes documents, and the
implementation of the retrieval, agent and knowledge-graph subsystems. It also
covers how data, storage and security were realised. The aim is to show the main
modules and how they fit together, rather than to reproduce the code in full.

\section{Technology Stack}

The platform was built from established, well-supported technologies, chosen so
that the team could concentrate on the design rather than on infrastructure.
Table~\ref{tab:stack} summarises the stack and the role of each part.

\begin{table}[H]
  \centering
  \caption{The technology stack and the role of each component. The user-facing
  application runs at the edge on Cloudflare, the processing and reasoning run on a
  GPU server, and the data lives in a managed database and object store, with
  language and vision models reached over the network.}
  \label{tab:stack}
  \small
  \begin{tabularx}{\textwidth}{@{}p{0.28\textwidth}X@{}}
    \toprule
    \textbf{Layer} & \textbf{Technology and role} \\
    \midrule
    Web application      & Next.js and React, deployed to a Cloudflare Worker through OpenNext, providing the user interface at the edge. \\
    Backend service      & FastAPI in Python, running on a GPU virtual machine, exposing the chat, agent, graph and pipeline interfaces. \\
    Database             & Supabase (PostgreSQL) with the pgvector extension for vector search and full-text search for keywords. \\
    Object storage       & Cloudflare R2 for the original files, the processed text, and the generated artefacts. \\
    Local models         & The BGE encoder for embeddings and a cross-encoder reranker, run on the server's GPU. \\
    Hosted models        & OpenAI and OpenRouter, providing the language models and the Qwen2.5-VL vision-language model. \\
    Document analysis    & DocLayout-YOLO for page layout and an optical character recognition toolkit for scanned text. \\
    Real-time            & Supabase Realtime, which delivers live chat and processing updates to the browser. \\
    \bottomrule
  \end{tabularx}
\end{table}

\section{The Ingestion Pipeline}

The ingestion pipeline turns a connected source into a searchable library. It is
implemented as a sequence of queue-driven stages, shown in
Figure~\ref{fig:pipeline}, each carried out by its own worker. A stage runs in
batches and only begins once the previous stage has completed for a document,
which makes the whole pipeline resumable: if a stage fails, it can be retried
without redoing the work that came before it.

\begin{figure}[H]
  \centering
  \includegraphics[width=\textwidth]{fig_pipeline.png}
  \caption{The implemented ingestion pipeline. Each stage is a separate worker
  that runs in batches and waits for the previous stage to finish, so processing
  is resumable and a failure in one document or stage does not lose the rest. The
  text and visuals are stored as artefacts, and the chunk embeddings are written
  to the vector database for retrieval.}
  \label{fig:pipeline}
\end{figure}

The stages map onto the worker modules of the backend. A sync worker fetches the
files from the connected source. A layout worker detects the regions of each page
using DocLayout-YOLO. An extraction worker reads the text, taking born-digital
text directly and reading scanned pages with the vision-language model through a
dedicated client. A captioning worker describes figures and tables. A chunking
worker splits the text into passages, an embedding worker encodes each passage
with the BGE model and writes it to the vector database, and a clustering worker
groups the passages into topics for coverage. A batch creator schedules the work
and a watchdog recovers stages that stall, so that large libraries are processed
steadily and reliably.

\section{Retrieval and the Grounded Chat}

The grounded chat is implemented around a runtime that carries out the
grade-then-answer flow from Chapter~4. When a question arrives, the chat interface
hands it to the runtime, which embeds the question, retrieves passages by vector
and by keyword, fuses the two ranked lists with reciprocal rank fusion, and
re-ranks the result with the cross-encoder. The agents that plan the retrieval,
grade the evidence and compose the answer are kept as separate, small steps, which
made each one easy to test and adjust. A retriever worker and a queue allow the
heavier retrieval work to run without blocking the interface.

The decision to answer or abstain is made explicitly after the evidence has been
graded. Listing~\ref{lst:abstain} shows this decision in simplified form: the
system only composes an answer when the graded evidence clears a confidence
threshold, and otherwise returns an honest abstention. The actual implementation
adds citation tracking and the option to refine the query and retrieve again
before abstaining.

\begin{lstlisting}[language=Python,caption={Simplified grade-then-answer decision in the chat runtime.},label={lst:abstain}]
graded = grade_evidence(question, reranked_passages)

if graded.is_sufficient and graded.confidence >= ABSTAIN_THRESHOLD:
    answer = compose_answer(question, graded.passages)
    return answer.with_citations(graded.passages)
else:
    return abstain("The selected library does not support an answer.")
\end{lstlisting}

\section{The Agent}

The agent is implemented as a small set of steps coordinated by an agent
interface. A language-model client targets a capable hosted model and returns
structured results, which are validated before use. The planning, data-extraction
and specification steps are separate functions, mirroring the design in
Chapter~4. Structured files are read exactly by a data module that preserves the
native cell types of a spreadsheet, while unstructured needs are met by reuse of
the retrieval stack. A specification module validates a chart or diagram and binds
the data values to it in code, and separate modules handle the rendering of
documents and images. Keeping the model's role to describing artefacts, and
binding data outside the model, is what makes the figures faithful.

\section{The Knowledge Graph}

The knowledge graph is implemented in four parts that follow its construction
flow. An extraction module prompts a language model to list the entities and
relationships in each chunk. A resolution module merges entities that refer to the
same thing under different names. A build module assembles the nodes and edges and
scores their weights, and an interface module serves the graph to the web
application, where it is drawn as an interactive diagram. Throughout, each node
and edge records the chunk it came from, so the graph remains connected to its
evidence.

\section{Data, Storage and Security}

The data described in Chapter~4 is stored in PostgreSQL through Supabase, with the
pgvector extension providing the vector search used by retrieval and a full-text
index providing the keyword search. The original files, the processed text and the
generated artefacts are kept in Cloudflare R2. Tenant isolation is enforced by
row-level security: every table that holds tenant data carries an organisation
identifier, and database policies ensure that a member can read only the rows that
belong to their organisation. This places the privacy guarantee at the database,
where it cannot be bypassed by a mistake in the application. Live updates, such as
new chat messages and changes in processing status, are delivered to the browser
through Supabase Realtime.

\section{Summary}

This chapter described the implementation of the platform. It set out the
technology stack, explained the queue-driven ingestion pipeline and its worker
modules, and described how the grounded chat, the agent and the knowledge graph
were built from small, separate steps that follow their designs. It also explained
how the data model, storage and the row-level security that guarantees tenant
isolation were realised. The next chapter presents the finished product through
the web application.
